% ============================================================================
% 3 PROBLEM STATEMENT AND RESEARCH QUESTIONS
% ============================================================================
\section{Problem Statement and Research Questions}
\label{sec:problem}

An organization can obtain plausible extraction and classification output from
a general LLM without training a task-specific model. Plausibility is not an
accuracy measure, however, and the organization usually lacks reference labels
for exactly the documents it wants to process. It must therefore solve two
linked engineering problems: evaluate a proposed method on representative
data, and create trustworthy annotations without making experts label every
document from scratch. Long documents add operational concerns---call volume,
partial failures, offset reconstruction and resumability---that short benchmark
sentences conceal.

We address the following questions:

\begin{description}[leftmargin=1.7cm, style=nextline, font=\bfseries, itemsep=1pt]
    \item[RQ1] How accurately and reliably does a training-free LLM reproduce
        reference labels for ticket classification and span extraction, under
        strict structural and semantic validation?
    \item[RQ2] How well does schema-guided pre-annotation localize Method,
        Task, Metric and Dataset mentions in coherent SciREX text using
        sentence-local predictions, and how does performance change from short
        to very long document windows?
    \item[RQ3] Can one reusable human-in-the-loop product turn model
        predictions into auditable reviewed data across scientific extraction
        and internal-ticket use cases?
\end{description}

The study is intentionally not a model leaderboard. A single hosted deployment
is fixed so that changes in data, prompt and pipeline remain visible. We use
strict span $F_1$ as the primary extraction measure, because a gold export must
reproduce exact offsets, and predeclare boundary-tolerant and overlap metrics as
secondary measures of whether a suggestion still directs a reviewer to the
useful phrase. For document-scale evaluation, source paper rather than window is
the statistical unit because several windows may originate from one paper.

Three expectations guide the analysis. First, syntactically valid JSON will not
guarantee correct labels or boundaries. Second, schema definitions and explicit
negative guidance will materially affect extraction quality. Third, a complete
document-scale run should be technically possible through sentence-level
checkpointing, but accuracy need not improve with longer windows because calls
remain sentence-local. These are evaluated only where retained artifacts support
the comparison; unexecuted configurations are not presented as results.
